\section{行为 - 不作为犯的成立条件}
\begin{enumerate}
  \item \zd{应为}。考的最多。要有作为义务。义务不能无中生有，来源根据：
        \begin{enumerate}
          \item 形式的四分说。传统观点，没有揭示作为义务真正的产生根据。
          \item \bt{实质二分说}。危险源 \yjt 危险流 \yjt| 法益对象。作为义务指要\zd{阻断消除危险源}，保护法益对象，虽然不是你创造的。
                \begin{enumerate}
                  \item 对危险源有\zd{监管义务}。
                        \begin{enumerate}
                          \item 危险物。\al{遛狗不牵狗绳、井盖}。
                          \item 他人的危险行为。行为人跟他人有关系，如监护、监管关系。
                          \item 先行行为。对法益创设了\zd{现实的危险}，不要混淆道德义务和刑法义务。
                                \begin{itemize}
                                  \item 违法阻却事由能成为 \blh 产生义务：
                                        \al{盗贼 \yjt
                                        被狗咬死 \yjt
                                        主人放任不管}（防卫过当，\zd{有防止过当结果发生的义务}）。
                                  \item 犯罪行为能成为 \blh 产生义务：\al{陕西盗墓土炕倒塌不救，导致夫妻摔死}。
                                \end{itemize}
                        \end{enumerate}
                  \item 对法益对象有\zd{保护义务}。
                        \begin{enumerate}
                          \item 特定关系。
                                \begin{itemize}
                                  \item 法律规定。\al{老妈和女友落水；老妈和老婆落水} \yjt \scx{义务冲突}（考纲删除）。
                                  \item 职务业务规定。消防员、警察。
                                  \item \zd{自愿救助行为}：\al{弃婴抱回家后有义务抚养送医等}。产生\zd{依赖关系}（排除他人救助可能）。
                                \end{itemize}
                          \item 特性领域。管理人 + 排他的支配作用（产生了依赖关系）。
                                \begin{itemize}
                                  \item \al{客人在家心脏病发作，主人有送医义务}。
                                  \item \al{分手我就从你家里跳下去}。
                                  \item \zd{领域扩大解释}，人的身体等都算。\al{大叔任由小女孩摸自己裸体}（不作为的猥亵儿童）。
                                \end{itemize}
                        \end{enumerate}
                \end{enumerate}
        \end{enumerate}
  \item 能为。要有\zd{履行能力}，作为的\zd{可能性}。
        \begin{enumerate}
          \item \al{孩子落水父亲腿抽筋} \yjt 已无作为可能性。
          \item \al{深圳强奸案，丈夫不敢去报警} \yjt 有作为可能性，片面的不作为的帮助犯。
        \end{enumerate}
  \item 不为。\zd{结果避免的可能性}。考过两次。
        \begin{enumerate}
          \item \al{甲车间操作不当乙必死无疑} \yjt 结果无法避免，\zd{履行已无意义}。
          \item \zd{刑法创设义务的意义是如果履行了义务，这个危害结果可能避免发生}。
          \item 所以死亡的因不能算到不作为救助上，而是过失操作上。
        \end{enumerate}
  \item 具有等价性：有\zd{程度要求}，要达到严重程度。\al{母亲一小时不喂养儿子和一个月不喂养}是有区别的。
  \item 不作为可以是过失 \hz 故意。有个和故意中\zd{认识错误}结合的考点：
        \begin{enumerate}
          \item 事实认知错误：\al{孩子落水看热闹，死的是自己孩子}。有看清是自己孩子的可能性，过失的不作为。
          \item 法律认识错误：\al{养子落水认为没有救助义务}。\zd{不知法者不免责}，看到还不去救，故意的不作为。
        \end{enumerate}
  \item 故意的不作为也有既遂未遂问题。考过一次。\al{把孩子遗弃到山洞；马路边；森林}。
  \item \zd{多个作为义务人之间存在先后，若第一作为义务人不救助 \hz 侵害法益，不能要求第二作为义务人实施救助}。
\end{enumerate}
